%!TeX root=../thesis.tex

\chapter{Conclusion}
\label{ch:conclusion}
\label{chap:conclusion}

This chapter closes the empirical and analytical work of the dissertation.
Section~\ref{sec:conclusion-rq} summarizes the main contributions in relation
to each research question. Section~\ref{sec:implications} discusses
implications for practitioners and researchers.
Section~\ref{sec:future-work} outlines the future work agenda that this study
opens.

% ─────────────────────────────────────────────────────
\section{Summary in Relation to the Research Questions}
\label{sec:conclusion-rq}

\paragraph{RQ1 --- How do software engineering teams use Hypo-Stage and ArchHypo
when working on a scaled product?}

The study produced a detailed observational account of tool adoption in two
teams over four weeks of real sprint work. The most consistent finding is
that hypothesis registration was senior-led: twelve of the thirteen registered
hypotheses were authored by staff or principal engineers who held explicit
architectural responsibilities. Junior and mid-level engineers were aware of
the tool and trained on the technique, but did not register hypotheses
independently within the observation window. Registration activity was
concentrated in two bursts --- one coinciding with the training session, and a
second when the Platform Team used the tool to structure a live architectural
discussion on Kafka failure-handling strategy --- rather than being continuously
distributed across sprints.

The integration of Hypo-Stage with the Backstage Internal Developer Portal
was consistently cited as an enabler of adoption. The
\texttt{EntityHypothesesTab} --- which surfaces only hypotheses linked to the
service a developer is already viewing --- reduced the cognitive cost of
context-switching into hypothesis management and was specifically mentioned as
a meaningful design improvement over a standalone tool. However, embedding
the tool into team-visible ceremonies (retrospectives, planning sessions)
did not occur spontaneously within the observation window; both teams noted that
this would require an explicit cultural agreement.

\paragraph{RQ2 --- What types and recurring structural shapes of architectural
uncertainties do teams identify, record, and manage using Hypo-Stage?}

The study produced thirteen architectural hypotheses across two teams. Analysis
of their source types, quality attributes, uncertainty and impact profiles, and
technical plan strategies characterizes the landscape of architectural
uncertainty that Hypo-Stage users made explicit (Section~\ref{sec:rq2-results}).
Requirements-sourced hypotheses dominated the Product Team (five of six), while
Solution-sourced hypotheses dominated the Platform Team's Kafka cluster (four of
five). The dominant quality attributes --- Reliability and Auditability ---
reflect the teams' real operational context: a scaled financial services product
operating under both internal security audit and privacy regulation obligations.

Beyond this characterization, five recurring structural archetypes were
identified from the data through thematic analysis: \textit{Hypothesis Cluster},
\textit{Regulatory Trigger Hypothesis}, \textit{Decoupling Seam Hypothesis},
\textit{Load-Validated Hypothesis}, and \textit{Shared Resource Governance
Hypothesis} (Section~\ref{sec:group-b-patterns}). These archetypes constitute the first empirically
grounded answer to the question of what shapes architectural hypotheses take when
ArchHypo is applied in a real industrial context, and they directly address a gap
identified as future work in \citet{silva2024patterns}. Readers should note
that these are \emph{emerging patterns}, derived from a thirteen-hypothesis
dataset; they require replication in additional organizational contexts before
they can be treated as consolidated contributions to the ArchHypo pattern
language (see Section~\ref{sec:future-work}).

\paragraph{RQ3 --- What benefits and challenges do team members perceive in
using Hypo-Stage?}

Participants across both teams reported that making architectural uncertainties
explicit and visible through Hypo-Stage had a positive effect on the quality of
their architectural discussions. The most salient benefit was the formalization
of tacit knowledge: by writing a hypothesis, engineers converted an informal
concern --- shared informally in chat or resolved verbally in a meeting --- into a
traceable, prioritized artifact with an explicit rationale. The Backstage
catalog integration was perceived as reducing cognitive overhead.

The dominant challenge remained the learning curve of the ArchHypo technique
itself. All four participants in the active listening sessions identified the
falsifiability requirement of hypothesis statements as non-intuitive under
delivery pressure, replicating findings from the PropSEP/Jetsoft evaluation
\citep{silva2025archhypo}. Sprint delivery pressure was named by both technical
leaders as the principal constraint on registration activity: formulating a
well-structured hypothesis requires focused, uninterrupted time that is scarce
in a product sprint. The team-behavior patterns in Group~A (Section~\ref{sec:group-a-patterns})
condense these findings into four actionable organizational conditions ---
\textit{Senior-Curated Hypothesis Backlog} (A1), \textit{Fast-Resolution Blind Spot} (A2),
\textit{Registration as a Slack-Dependent Practice} (A3), and
\textit{Negative-Formulation Barrier} (A4) --- that jointly explain why adoption is
sustainable only when specific organizational enablers are in place. Like the
Group~B archetypes, these are emerging propositions and require further
validation.

\paragraph{RQ4 --- Which refinements to Hypo-Stage and its usage guidelines
emerge from the empirical study?}

Nine categories of empirically grounded refinements were implemented during the
observation period and documented in Section~\ref{sec:refinements}, each traced to a concrete
observation and a commit in the Hypo-Stage repository. The refinements span
four layers of the tool: Backstage integration and catalog surfacing (making
hypotheses discoverable without navigation); backend API and filtering (enabling
team-scoped and status-filtered views); the prioritization dashboard
(introducing focus panels and urgency signals derived from uncertainty and impact
combinations); and form usability (adding inline ArchHypo guidance to reduce the
vocabulary barrier at the moment of registration).

The most structurally significant refinement is the evolution chart introduced
in Section~\ref{subsec:refinement-chart} (Figure~\ref{fig:evolution-chart}), which plots the shift in a
hypothesis's uncertainty and impact ratings over successive assessments. This
feature directly addresses the absence of assessment history in the initial
Hypo-Stage version \citep{correa2025} and enables teams to see whether their
technical plan actions are having the intended effect of reducing uncertainty
or impact. Together, the nine refinement categories operationalize a cycle of
\emph{observe $\rightarrow$ refine $\rightarrow$ redeploy $\rightarrow$ observe} that is consistent with the
action-research element of the research design.

% ─────────────────────────────────────────────────────
\section{Implications for Practice and Research}
\label{sec:implications}

\paragraph{For practitioners adopting hypothesis-engineering tools.}

The study demonstrates that tool integration within the developer's daily
environment is a meaningful enabler of adoption. Deploying Hypo-Stage as a
Backstage plugin --- where engineers already manage services, documentation, and
CI/CD pipelines --- removed a class of adoption friction that a standalone tool
would have created. This suggests a general principle: hypothesis-engineering
tools are more likely to be adopted when they are delivered as extensions to
an existing workflow rather than as a new workflow to be added.

The finding that registration remains senior-led, even with tool support, has
practical implications for organizations intending to distribute architectural
decision-making. Tool availability is a necessary but not sufficient condition
for broad participation. Organizations should combine tool deployment with
explicit protected time for architectural hygiene (see Group~A pattern A3) and
with statement scaffolding that lowers the cognitive barrier to formulating
falsifiable hypotheses (see Group~A pattern A4 and Hypo-Stage form improvements
in Section~\ref{subsec:refinement-forms}). The Group~B archetypes can serve immediately as a
recognition vocabulary for experienced hypothesis engineers encountering
familiar situations in new projects.

\paragraph{For researchers studying tool-mediated architectural practices.}

This study is the first evaluation of ArchHypo in a tool-mediated setting and
the first to produce a dataset of architectural hypotheses large enough to
support structural classification. Several methodological observations may
benefit future work. First, the formative evaluation design --- in which
refinements are implemented during the observation period --- provides richer
insight into \emph{why} difficulties arise and \emph{which} changes address
them than a retrospective evaluation would. This design is well-suited to tool
artifacts whose value proposition is tightly coupled to usability and workflow
fit. Second, the dataset of thirteen hypotheses, while small, demonstrates that
even a short observation window in a naturalistic setting produces enough
artifacts to support initial pattern identification. Larger datasets will be
needed to test frequency and generalizability.

The observation that hypothesis registration is senior-led is consistent with
findings from both prior ArchHypo evaluations \citep{silva2024ieeesw,
silva2024tse} and with the broader literature on tacit knowledge in architectural
decision-making \citep{souza2019}. This convergence across different settings,
time periods, and methods provides some basis for treating senior-anchor
dependency as a robust characteristic of early ArchHypo adoption rather than an
artifact of any single study.

% ─────────────────────────────────────────────────────
\section{Limitations and Future Work}
\label{sec:future-work}

\subsection*{Study Limitations}

\paragraph{Single organizational context and Backstage dependency.}
The study was conducted within one technology organization that already operated
a mature Backstage deployment. The enabling condition of Backstage availability
means that findings about adoption friction may not transfer directly to
organizations without an Internal Developer Portal. The organizational culture,
team size, and product complexity of the study site may differ substantially from
other contexts in which researchers or practitioners seek to apply these findings.

\paragraph{Small hypothesis dataset and short observation window.}
Thirteen hypotheses produced over four weeks by two teams is a limited dataset.
The Group~A and Group~B patterns are derived from this dataset and must be treated
as emerging propositions rather than consolidated findings. The four-week window
is also too short to observe full hypothesis lifecycle completion (validation or
invalidation); eleven of the thirteen hypotheses remained in Open status at the
end of the period. A longer observation window would be needed to study how
teams manage hypotheses through to resolution and whether the tool adequately
supports retrospective learning after a hypothesis closes.

\paragraph{Researcher dual role.}
The researcher's dual role as organizational employee and academic investigator
is a known source of potential bias. Although strategies were employed to
maintain role separation (Section~\ref{sec:researcher-role}), the researcher's presence as a tool
author and subject-matter expert may have influenced how participants described
their experience of the technique, introducing social desirability effects.

\paragraph{No comparison condition.}
The study does not include a control group or comparison condition against teams
using ArchHypo without Hypo-Stage, or teams not using any structured
uncertainty-management approach. This limits the ability to attribute observed
benefits specifically to the tool rather than to the introduction of a structured
technique in general.

\subsection*{Future Work Agenda}

\paragraph{Replication and longitudinal extension.}
The most immediate priority is replication of this study in additional
organizational contexts --- including organizations without an existing Backstage
deployment --- to test whether the adoption patterns, tool refinements, and
emerging archetypes are context-specific or more broadly applicable.
A longitudinal design spanning multiple development cycles would allow the
study of hypothesis lifecycle completion, assessment evolution, and whether
teams internalize the technique over time to the point where senior-anchor
dependency diminishes.

\paragraph{Consolidating the Group~B archetypes into the ArchHypo pattern language.}
The five technical architectural hypothesis archetypes (B1--B5) were identified
from a single dataset and are explicitly preliminary. Consolidating them
requires: (a)~attempting to identify or refute each archetype in new hypothesis
datasets from different organizations and domains; (b)~subjecting each pattern
to the structured review process used for the existing ArchHypo pattern language
\citep{silva2024patterns}; and (c)~expanding the catalog to cover quality
attributes not represented in this study (e.g., security-only hypotheses,
performance-at-scale hypotheses in non-financial contexts).

\paragraph{Validating the Group~A team-behavior patterns.}
The four organizational conditions in Group~A (A1--A4) are equally preliminary.
Future work should design intervention studies that test whether addressing
these conditions --- for instance, by introducing protected architectural hygiene
sessions (A3) or embedding statement templates in the hypothesis form (A4) ---
measurably increases registration breadth beyond senior engineers, reduces
the formulation time reported by participants, or improves hypothesis quality
as judged by the teams themselves.

\paragraph{Broadening participation beyond senior engineers.}
A central open question is whether and how hypothesis registration can become a
practice accessible to engineers at all levels, not only to those who already
hold explicit architectural responsibilities. Research directions include:
guided hypothesis templates that lower the formulation barrier (already initiated
in the form usability refinements); asynchronous architectural review workflows
that do not require simultaneous availability of a senior engineer; and
gamified or social nudge mechanisms that encourage mid-level engineers to
surface uncertainties they currently resolve tacitly.

\paragraph{Integration with Architecture Decision Records.}
Hypo-Stage manages the pre-decision lifecycle of architectural uncertainty.
Once a hypothesis is validated or invalidated, the resulting decision is
currently not captured in any structured artifact within the tool. Future
work should explore the integration of Hypo-Stage with ADR tools such as
Log4brains, so that the evidence chain from hypothesis registration through
technical plan execution to final architectural decision is preserved as a
single traceable thread.

\paragraph{Quantitative assessment of tool impact.}
This study does not include any quantitative measure of the impact of
Hypo-Stage on decision quality, architecture drift, or delivery velocity.
Future work using a quasi-experimental or interrupted time-series design could
compare architectural decision quality or technical debt accumulation between
periods with and without active hypothesis management, providing a more
rigorous basis for claims about the tool's organizational value.

\paragraph{Extending Hypo-Stage to multi-team hypothesis governance.}
The present evaluation involved two co-located teams within the same
organization. As organizations scale, hypotheses registered by one team
frequently have cross-team implications --- this was already visible in the blob
governance hypotheses (H07, H08) and the platform Kafka cluster (H09--H13).
Future work should explore cross-team hypothesis visibility, shared ownership
models, and notification mechanisms that make cross-boundary hypotheses
discoverable to the teams they affect, without creating coordination overhead
that undermines team autonomy.
